% !TeX root = ../../Book1.tex
% 纲要标题：### 9.4 结构定理的计算与模论解释
% 纲要位置：计划纲要.md，第 383 行。

\section{结构定理的计算与模论解释}
\label{b1:sec:0904}

结构定理给出了分类形式，实际计算还须把生成元、关系和坐标变换一一对应。下面先固定关系矩阵的行列约定，再用一个完整例子说明怎样从化简结果读出循环因子并追踪具体元素。随后讨论整数格的商群、子格指数与子式判据。节末的模论解释只是对第二卷的预告，本节全部计算仍只依赖整数矩阵、子群与商群。

% 纲要内容（仅作写作计划，尚非教材正文）：
%
% * 由整数关系矩阵确定交换群
% * 整数格的商群与有限指数
% * 模论解释与后续推广预告
%
% ---
%

\subsection{整数关系矩阵与交换群}
\label{b1:subsec:0904:01}

使用结构定理计算时，第一步是固定矩阵的方向。本节一律把生成元列为坐标轴，把一条整数关系写成一个列向量。给定 \(m\times n\) 整数矩阵 \(M\)，可以用它的各列施加关系，并定义交换群
\[
A_M\coloneqq\mathbb Z^m/M\mathbb Z^n.
\]
这里 \(M\mathbb Z^n\) 是由这些列关系的全部整数线性组合组成的子群。列操作改变关系生成集，行操作改变群的生成坐标；两者都必须是可逆的整数操作。

\ProseParagraphBreak
若要用已知关系分析一个事先给定的交换群 \(A\)，还须核对关系是否完整。设 \(A\) 由 \(a_1,\ldots,a_m\) 生成，并令
\[
\pi:\mathbb Z^m\longrightarrow A,
\qquad e_i\longmapsto a_i
\]
为自然满同态。若 \(M\) 的各列只记录若干已经验证的关系，则只能得到 \(M\mathbb Z^n\subseteq\ker\pi\)，因而 \(\pi\) 诱导自然满同态
\[
\overline\pi:\mathbb Z^m/M\mathbb Z^n\longrightarrow A.
\]
只有进一步证明 \(M\mathbb Z^n=\ker\pi\)，这个满同态才是同构，Smith 化简所得结构才是 \(A\) 本身的分类。例如 \(A=C_2=\langle a\rangle\) 确实满足关系 \(4a=0\)，但只记录这一关系得到的是 \(\mathbb Z/4\mathbb Z\cong C_4\)，它只自然满射到 \(C_2\)；遗漏的关系 \(2a=0\) 正说明所列关系并不完整。

\ProseParagraphBreak
Smith 计算本身的输入是一张有限整数矩阵 \(M\)。在 \(\mathbb Z\) 上逐次作可逆行列操作，并同时累计这些操作，输出 \(m\times m\) 与 \(n\times n\) 整数可逆矩阵 \(U,V\)，以及 Smith 标准形 \(D=UMV\)。只求同构类型时，从 \(D\) 读取自由秩与非平凡不变因子即可；还要追踪原生成元或具体元素时，则须保留 \(U,V\) 及其逆矩阵。\cref{b1:lem:smith-integer}的证明给出了这一过程的终止性，但本章不讨论它的复杂度。若 \(M\) 用来分析一个事先给定的群，求得生成完整关系核的矩阵是 Smith 化简之前另须解决的问题。

\ProseParagraphBreak
设 Smith 标准形有 \(t\) 个非零对角元 \(d_1\mid\cdots\mid d_t\)。\cref{b1:tab:smith-coordinate-types}区分新坐标中的三种情况：前 \(t\) 个坐标受到非零关系约束，其余 \(m-t\) 个坐标没有这样的约束。

\begin{table}[htbp]
\centering
\caption{Smith 标准形中的坐标与商群因子}
\label{b1:tab:smith-coordinate-types}
\begin{tabularx}{\linewidth}{@{}>{\raggedright\arraybackslash}p{0.44\linewidth}>{\raggedright\arraybackslash}X@{}}
\toprule
新坐标中的情形 & 对商群的贡献 \\
\midrule
非单位主元 \(d_i>1\) & 非平凡有限循环因子 \(C_{d_i}\) \\
单位主元 \(d_i=1\) & 该坐标在商群中为零，平凡因子 \(C_1\) 删去 \\
没有非零主元约束的坐标 & 每个贡献一个无限循环因子，共有 \(m-t\) 个 \\
\bottomrule
\end{tabularx}
\end{table}

矩阵的非零 Smith 对角元允许为 \(1\)；群结构定理所列的不变因子则删去这些 \(1\)，只保留大于 \(1\) 的项。标准形中的零列表示经过可逆整数重组后出现的零关系，可以在这一阶段删去；这不意味着原关系组中必有某条关系可以直接删除。关系矩阵可以有任意多个相互依赖的列；自由部分的秩由目标格的坐标数减去非零主元数得到，为 \(m-t\)，不是 \(n-t\)。

\ProseParagraphBreak
例如对一生成元群施加关系 \(2x=0\)、\(3x=0\)，关系矩阵为 \(M=\begin{pmatrix}2&3\end{pmatrix}\)。取整数可逆矩阵
\[
V=\begin{pmatrix}-1&3\\1&-2\end{pmatrix},
\qquad \det V=-1,
\]
便有
\[
MV=\begin{pmatrix}1&0\end{pmatrix}.
\]
零列是在两个原关系经过可逆重组后产生的；原关系中却没有一条可以直接删去，因为两条关系合起来定义平凡群，只保留 \(2x=0\) 或只保留 \(3x=0\) 时分别得到 \(C_2\) 或 \(C_3\)。

\paragraph*{求同构类型：读取 Smith 对角元}
例如群由 \(x,y\) 生成，关系为 \(4x+2y=0\)、\(6x+8y=0\)，矩阵为
\[
M=\begin{pmatrix}4&6\\2&8\end{pmatrix}.
\]
一次完整化简为
\[
\begin{pmatrix}4&6\\2&8\end{pmatrix}
\longrightarrow
\begin{pmatrix}2&8\\4&6\end{pmatrix}
\longrightarrow
\begin{pmatrix}2&8\\0&-10\end{pmatrix}
\longrightarrow
\begin{pmatrix}2&0\\0&-10\end{pmatrix}
\longrightarrow
\begin{pmatrix}2&0\\0&10\end{pmatrix}.
\]
四步分别为交换两行、第二行减第一行的两倍、第二列减第一列的四倍、第二行变号。因此 \(A\cong C_2\oplus C_{10}\)。

\paragraph*{追踪原元素：保存行变换矩阵}
若还需要指出商群元素怎样对应，不能只记住对角数字。本例行操作的总矩阵为
\[
U=\begin{pmatrix}0&1\\-1&2\end{pmatrix},\qquad
U^{-1}=\begin{pmatrix}2&-1\\1&0\end{pmatrix}.
\]
原坐标向新坐标转换使用 \(x\mapsto Ux\)；若要把新坐标轴对应的生成元写回原坐标，则使用 \(U^{-1}e_i\)。先在整数坐标上定义同态
\[
\psi:\mathbb Z^2\longrightarrow C_2\oplus C_{10},\qquad
\psi(u,v)=\Bigl([v]_2,[-u+2v]_{10}\Bigr).
\]
前述列操作的矩阵及其整数逆矩阵为
\[
V=\begin{pmatrix}1&-4\\0&1\end{pmatrix},\qquad
V^{-1}=\begin{pmatrix}1&4\\0&1\end{pmatrix}.
\]
实际相乘得到 \(UMV=\operatorname{diag}(2,10)\)。因为 \(V\mathbb Z^2=\mathbb Z^2\)，有 \(U(M\mathbb Z^2)=2\mathbb Z\oplus10\mathbb Z\)。因此
\[
\ker\psi=U^{-1}(2\mathbb Z\oplus10\mathbb Z)=M\mathbb Z^2.
\]
\(U\) 在整数上可逆，随后分别取模的映射满射，所以 \(\psi\) 满射。

\ProseParagraphBreak
由\cref{b1:thm:first-isomorphism}，\(\psi\) 诱导同构 \(\mathbb Z^2/M\mathbb Z^2\rightarrow C_2\oplus C_{10}\)，把陪集 \((u,v)+M\mathbb Z^2\) 送到 \(\psi(u,v)\)。原坐标中的 \(U^{-1}e_1=(2,1)\) 对应 \(C_2\) 的标准生成元，\(U^{-1}e_2=(-1,0)\) 对应 \(C_{10}\) 的标准生成元；它们在商群中的阶分别为二与十。

\ProseParagraphBreak
若生成元为 \(x,y,z\)，关系只有 \(2x=0\)、\(x+3y=0\)，则先用第二条消去 \(x=-3y\)，第一条成为 \(6y=0\)，而 \(z\) 没有关系。因此群为 \(C_6\oplus\mathbb Z\)。没有非零对角关系约束的位置贡献自由生成元，而不是有限阶循环因子。

\subsection{整数格的商群与有限指数}
\label{b1:subsec:0904:02}

\begin{proposition}[子格指数与商群结构]\label{b1:prop:lattice-index}
设 \(m,n\) 为正整数，\(M\) 为 \(m\times n\) 整数矩阵，\(L=M\mathbb Z^n\le\mathbb Z^m\)。若 \(M\) 的 Smith 标准形有 \(t\) 个非零对角元 \(d_1,\ldots,d_t\)，则
\[
\mathbb Z^m/L\cong\mathbb Z^{m-t}\oplus
\bigoplus_{i=1}^t C_{d_i}.
\]
特别地，\(L\) 的指数有限当且仅当 \(t=m\)，此时指数为 \(\prod_i d_i\)。若 \(M\) 为非奇异的 \(m\times m\) 矩阵，则指数为 \(|\det M|\)。
\end{proposition}

\begin{proof}
由\cref{b1:lem:smith-integer}，存在整数可逆矩阵 \(U,V\) 使 \(UMV=D\)，其中 \(D\) 是所述对角阵。由于 \(V\mathbb Z^n=\mathbb Z^n\)，有 \(U(L)=D\mathbb Z^n=\langle d_1e_1,\ldots,d_te_t\rangle\)。因此 \(x+L\mapsto Ux+U(L)\) 给出商群同构。再把前 \(t\) 个坐标分别模 \(d_i\)，其余坐标保留，所得满同态的核恰为 \(U(L)\)；\cref{b1:thm:first-isomorphism}遂给出所述商群公式。

指数就是商群的元素数。当 \(m-t>0\) 时有无限循环直和因子，故指数无限；当 \(m=t\) 时指数为有限循环因子阶的乘积。若 \(M\) 是非奇异方阵，则 \(t=m\)，而 \(\det U,\det V\in\{1,-1\}\)，所以从 \(UMV=D\) 得 \(|\det M|=\det D=\prod_i d_i\)。
\end{proof}

例如上一小节的二维关系子格指数为二十，等于 \(|4\cdot8-6\cdot2|\)。但行列式只给出总阶，并不决定商群：对角矩阵 \(\operatorname{diag}(1,20)\) 与 \(\operatorname{diag}(2,10)\) 的行列式同为二十，商群分别为 \(C_{20}\) 与 \(C_2\oplus C_{10}\)，两者中元素的最大阶不同。

\begin{proposition}[行列式因子]\label{b1:prop:determinantal-divisors}
设 \(m,n\) 为正整数，\(M\) 为 \(m\times n\) 整数矩阵。对 \(1\le k\le\min(m,n)\)，令 \(\Delta_k(M)\) 为 \(M\) 的全部 \(k\) 阶子式的非负最大公因数，并将它称为第 \(k\) 个\term[xinglieshi-yinzi]{行列式因子}[determinantal divisor]；全部相应子式都为零时约定其最大公因数为零，并令 \(\Delta_0(M)=1\)。令
\[
t=\operatorname{rank}_{\mathbb Q}M
\]
为 \(M\) 在有理数域上的矩阵秩。这个数也等于像子群 \(M\mathbb Z^n\) 的自由秩，以及 Smith 标准形中非零对角元的个数。将这些非零对角元依次记为 \(d_1,\ldots,d_t\)，则
\[
\begin{aligned}
\Delta_k(M)&=d_1\cdots d_k &&(1\le k\le t),\\
\Delta_k(M)&=0 &&(t<k\le\min(m,n)).
\end{aligned}
\]
因此比值公式为
\[
d_k=\frac{\Delta_k(M)}{\Delta_{k-1}(M)}\qquad(1\le k\le t).
\]
\end{proposition}

\symidx{\(\Delta_k(M)\)：矩阵的行列式因子}

\begin{proof}
整数可逆矩阵在有理数域上也可逆，所以 Smith 化简 \(UMV=D\) 保持 \(\operatorname{rank}_{\mathbb Q}\)，而 \(D\) 的有理秩就是其非零对角元个数。另一方面，\(U\) 把像子群 \(M\mathbb Z^n\) 同构地送到由 \(d_1e_1,\ldots,d_te_t\) 生成的自由交换群，所以该像子群的自由秩也等于 \(t\)。

整数初等行列操作保持各阶子式的最大公因数。以某行加另一行的整数倍为例，每个受影响子式由行列式的线性性写成原子式与另一原子式的整数线性组合；如果所加行已包含在子式中，额外项有重复行而为零。逆操作给出反向整除，所以最大公因数不变。换行、变号及列操作同样成立。

对 Smith 对角阵，当 \(1\le k\le t\) 时，非零 \(k\) 阶子式是所选 \(k\) 个对角元的乘积。若指标为 \(i_1<\cdots<i_k\)，则 \(i_j\ge j\)，整除链使 \(d_1\cdots d_k\) 整除该乘积；而前 \(k\) 个对角元的子式恰好等于此乘积。因此它就是全部子式的最大公因数。各项为正，结合 \(\Delta_0(M)=1\)，便得到所写的比值公式。

若 \(k>t\)，对角阵的任意 \(k\) 阶子式都为零，故 \(\Delta_k(M)=0\)。特别地，零矩阵的所有正阶行列式因子均为零，不适用上述比值公式。
\end{proof}

\begin{corollary}[满行秩矩形矩阵的子格指数]\label{b1:cor:full-row-rank-index}
设 \(m,n\) 为正整数，\(M\) 为满行秩的 \(m\times n\) 整数矩阵，即 \(\operatorname{rank}_{\mathbb Q}M=m\)。则 \(m\le n\)，并且
\[
[\mathbb Z^m:M\mathbb Z^n]=\Delta_m(M).
\]
即子格 \(M\mathbb Z^n\) 在 \(\mathbb Z^m\) 中的指数等于全部 \(m\) 阶子式绝对值的最大公因数。若 \(m=n\)，这就是 \(|\det M|\)。
\end{corollary}

\begin{proof}
满行秩意味着非零 Smith 对角元的个数 \(t=m\)。又因 \(t\leq\min(m,n)\)，必有 \(m\leq n\)。由\cref{b1:prop:lattice-index}，指数为 \(d_1\cdots d_m\)；再由\cref{b1:prop:determinantal-divisors}，此乘积等于 \(\Delta_m(M)\)。方阵只有一个 \(m\) 阶子式，即其行列式。
\end{proof}

这里必须是满行秩，而不能笼统地只说矩形矩阵“满秩”。若 \(n<m\) 且 \(M\) 满列秩，则 \(t=n<m\)，商群仍有自由秩 \(m-n\)，因而子格指数无限；此时全部最大非零子式是 \(n\) 阶子式，它们的最大公因数并不表示 \(M\mathbb Z^n\) 在 \(\mathbb Z^m\) 中的有限指数。

\ProseParagraphBreak
前面用于说明零列的矩阵 \(M=\begin{pmatrix}2&3\end{pmatrix}\) 有满行秩，而
\[
M\mathbb Z^2=2\mathbb Z+3\mathbb Z=\mathbb Z,
\]
因为 \(3-2=1\)。其指数为 \(1=\gcd(2,3)\)，不能任选一个非零最大阶子式 \(2\) 或 \(3\) 作为指数。

\ProseParagraphBreak
行列式因子也适合核验小矩阵的 Smith 结果。本节的二维算例中，全部元素的最大公因数为二，唯一二阶子式为二十，因此 \(d_1=2,d_2=10\)。大矩阵的子式很多，可以先做整数行列化简，再抽查少量子式以发现计算失误。不过，若所抽取子式的最大公因数为 \(g\)，只能直接得到 \(\Delta_k(M)\mid g\)，不能据此断言相等。例如 \(\operatorname{diag}(2,3)\) 的全部一阶子式最大公因数为 \(1\)，只检查条目 \(2\) 就不能恢复它。

\ProseParagraphBreak
若要独立确认正的候选值 \(\delta=\Delta_k(M)\)，须证明全部 \(k\) 阶子式均被 \(\delta\) 整除，并找到最大公因数恰为 \(\delta\) 的一组子式。这两个方向合起来才给出等号。另一种充分验证是完整保存 \(U,V\)，核验 \(UMV=D\) 及 \(U,V\) 的整数逆矩阵，并检查 \(D\) 的非零对角元为正且形成整除链；这样便确认了 Smith 结果，无须枚举全部子式。

\subsection{模论解释与后续推广预告}
\label{b1:subsec:0904:03}

本章始终只用整数倍、子群与商群。下一卷将把这种结构重新解释为模：整数 \(n\) 对交换群元素 \(x\) 的作用就是 \(nx\)。这个预告不作为当前证明的输入；模的定义与公理将在第二卷第二章“模、子模与商模”正式建立。

\ProseParagraphBreak
从这一后续观点看，\(\mathbb Z^m\) 是有限秩自由对象，关系矩阵记录了从一个自由对象到另一个自由对象的映射，商群记录生成元在关系约束下留下的信息。整数的任意理想由一个元素生成，这一性质将被抽象为主理想整环条件，从而把本章的循环分解推广到更一般的系数系统。在本章，支持消元的具体工具是整数除法以及 Bézout 等式；推广时必须重新核对系数系统是否提供相应工具。

\ProseParagraphBreak
例如对有理数系数多项式也可以作带余除法，第二卷第四章“主理想整环上的模与标准形”会用类似方法研究线性变换的有理标准形。届时矩阵中的“整数关系”变为“多项式关系”，循环因子表达一个向量及其在反复作用下的生成规律。这个类比解释了有限生成交换群为什么安排在群论之后：它既完成一项独立分类，也提供了从群的生成关系过渡到模的生成关系的具体经验。

\ProseParagraphBreak
实际使用本章结果时，可按三步整理资料。第一步，确认群交换且有限生成，选择有限生成元，并取得一张列向量生成整个关系核的有限整数矩阵。第二步，明确行列约定，在整数上作可逆变换，同时累计 \(U,V\)，直至得到 \(D=UMV\)。第三步，核验 \(U,V\) 都有整数逆矩阵，且 \(D\) 的非零对角元为正并形成整除链，再从 \(D\) 读取自由秩及非平凡不变因子。

\ProseParagraphBreak
有限秩子群定理只保证完整关系核存在有限基，并不自动给出可供计算的生成组。若矩阵只记录部分已知关系，只能得到相应展示群到原群的自然满同态，不能把展示群的分类结果直接当作原群的分类。若需要追踪特定元素或子群，应使用保存的坐标变换；若只需判断两个有限生成交换群是否同构，比较自由秩与归一化后的不变因子即可。
